\documentclass[10pt,a4paper]{amsart}
\usepackage[margin=19mm]{geometry}
\usepackage{amsmath,amssymb,mathtools}
\usepackage[scr=rsfs,cal=euler]{mathalfa}
\usepackage{xfrac}
\usepackage[unicode,hidelinks,bookmarks=false]{hyperref}
\newcommand{\ie}{i.e.\ }
\newcommand{\cf}{cf.\ }
\newcommand{\resp}{respectively\ }
\newcommand{\Subsection}{\subsection}
\providecommand{\nobreakdash}{\nobreak\hbox{-}}
\setlength{\emergencystretch}{2em}
\multlinegap0pt
\usepackage{bm}
\usepackage{bbm}
\usepackage{dsfont}
\usepackage{xspace}
\usepackage{cancel}
\usepackage{mathtools}
\usepackage{amsthm}
\newtheorem{lemma}{Lemma}
\newcommand{\Cov}{\mathbb{C}\text{ov}}
\newcommand{\IdenMat}{\textbf{I}}
\newcommand{\LogCoe}{\boldsymbol{\gamma}}
\newcommand{\ObsCov}{\textbf{x}_i}
\newcommand{\ObsRes}{\textbf{y}_i}
\newcommand{\SumOverN}{\sum_{i=1}^n}
\newcommand{\vect}{\emph{vec}}
\title[Doubly robust estimation with functional outcomes missing at]{Doubly robust estimation with functional outcomes missing at random}
\author{Xijia Liu, Kreske Felix Ecker, Lina Schelin, Xavier de Luna}
\date{Source: arXiv:2411.17224v3 (CC BY 4.0)}
\begin{document}
\maketitle

\noindent\emph{Source and licence.} Doubly robust estimation with functional outcomes missing at random. Version arXiv:2411.17224v3 (CC BY 4.0), distributed under the Creative Commons Attribution 4.0 International licence, as recorded in the arXiv licence metadata for this exact version and shown on the abstract page https://arxiv.org/abs/2411.17224v3. Full rights notice: licenses/arxiv-2411.17224-CC-BY-4.0.txt.\par\smallskip

\noindent\textit{Editorial scope.} Counted unit: the Lemma whose source label is \texttt{lemma:finite\_DR} in the article (source0.tex lines 437--442, source label \texttt{lemma:finite\_DR}), delivered together with the article's own complete original proof of it (source0.tex lines 444--604, one \texttt{proof} environment of 161 lines). No mathematics is written, repaired or re-derived: statement and proof are the article's own text line for line. Required context retained uncounted: equ:regmod1 (lines 111--113); equ:ps (lines 116--118); equ:regmod2 (lines 422--424). Every definition, display and label of those lines is retained unchanged. Omitted from this package and not redistributed: 1--110, 114--115, 119--421, 425--436, 443--443, 605--808. Nothing inside a retained excerpt is omitted or altered: every retained body is delivered exactly as the article's lines are written, so no omission receipt of any kind accompanies this package. Citations: the retained excerpts carry no citation command at all, so no bibliography entry of the article is transcribed and no reference is redistributed. Reference adaptation: none was needed and none was made; every reference inside the delivered unit resolves inside it, so no unresolved reference or citation is produced. Printed slips kept literal: no printed word, symbol, hypothesis or number of the counted statement or of its proof was altered. Layout and class: the article's own class is \texttt{article} with packages including inputenc, mathrsfs, graphicx, amsfonts, amsmath, bm, bbm, amssymb, amsthm, xcolor, multirow, natbib, placeins, geometry; the wrapper is the lane's pinned \texttt{amsart} editorial wrapper at 10pt on A4, which restates the theorem-like environments the retained text needs and adds the article's own macro definitions that the retained text needs (\texttt{\textbackslash Cov}, \texttt{\textbackslash IdenMat}, \texttt{\textbackslash LogCoe}, \texttt{\textbackslash ObsCov}, \texttt{\textbackslash ObsRes}, \texttt{\textbackslash SumOverN}, \texttt{\textbackslash vect}), with extra wrapper packages (bm, bbm, dsfont, xspace, cancel, mathtools). The delivered numbering is the unit's own. Assets: no retained span contains a figure environment, a graphics-inclusion command, a PDF-page inclusion, a table or a TikZ picture; no image file is copied into this package, the package declares no asset dependency and no image file is redistributed with it. Licence: version arXiv:2411.17224v3 of this article is distributed under the Creative Commons Attribution 4.0 International licence, as recorded in the arXiv licence metadata for this exact version and shown on the abstract page https://arxiv.org/abs/2411.17224v3. A full rights notice is delivered at licenses/arxiv-2411.17224-CC-BY-4.0.txt. Characters: every retained excerpt is pure ASCII, so the delivered engine, XeLaTeX with its default Latin Modern text font, reports no missing character.
        \begin{equation}\label{equ:regmod1}
            \mathcal{Y}_i = \ObsCov^{\top}\boldsymbol{\beta} + \epsilon_i,
        \end{equation}

            \begin{equation}\label{equ:ps}
               \tau(\ObsCov^{\top} \LogCoe),
            \end{equation}

        \begin{equation}\label{equ:regmod2}
            \ObsRes = \textbf{B}^{\top} \ObsCov + \boldsymbol{\epsilon}_i
        \end{equation}

    \begin{lemma}\label{lemma:finite_DR}
       Assume that at least one of the two working models \eqref{equ:regmod1} and \eqref{equ:ps} is correctly specified, i.e., $E(\ObsRes \mid \ObsCov)= \textbf{B}^{\top} \ObsCov$ or $\Pr[ Z_i=1 | \ObsCov ]= \tau(\ObsCov^{\top} \LogCoe)$. Then,  $\sqrt{n}(\hat{\boldsymbol{\mu}}_{DR} - \boldsymbol{\mu}_y)$ has an asymptotic multivariate normal distribution with mean zero. In the case that both working models are correctly specified, the covariance matrix simplifies and we have:
       \begin{align*}
           \sqrt{n}(\hat{\boldsymbol{\mu}}_{DR} - \boldsymbol{\mu}_y) \overset{d}{\rightarrow} \mathcal{N}\left(\textbf{0}, \textbf{B}^{\top} \boldsymbol{\Sigma}_{x} \textbf{B} + \mathbb{E} [\tau^{-1}(\ObsCov^{\top}\LogCoe)] \boldsymbol{\Sigma}_{\epsilon}\right).
       \end{align*} 
    \end{lemma}

    \begin{proof}       
        Based on the two models, the doubly robust estimator of the expected value of the outcome variables, $\boldsymbol{\mu}_{y} = \mathbb{E}(\textbf{y})$, is defined as
            \begin{equation*}
                \hat{\boldsymbol{\mu}}_{DR} = n^{-1} \SumOverN \left( \hat{\textbf{B}}^{\top}\ObsCov + \frac{Z_i}{\tau(\ObsCov^{\top}\hat{\LogCoe})}(\ObsRes -\hat{\textbf{B}}^{\top}\ObsCov)  \right),
            \end{equation*}
        where $\hat{\textbf{B}}$ and $\hat{\LogCoe}$ are maximum likelihood estimators.
        We apply the M-estimation approach to derive the asymptotic distribution of the doubly robust estimator.
        To do so, we first define the vector equations $G_n(\boldsymbol{\theta}) = n^{-1}\SumOverN{\psi_i(D_i, \boldsymbol{\theta})}$, where $D_i$ denotes the data of $i$th observation, and $\boldsymbol{\theta} = (\boldsymbol{\mu}_y, \vect(\textbf{B}), \LogCoe)^{\top}$ denotes the vector of all parameters. 
        To simplify the notations, we ignore $D_i$ in the equation $\psi_i$.
        The estimator $\hat{\boldsymbol{\theta}}$, which is the solution of the vector equations $G_n(\boldsymbol{\theta}) = 0$, has the following limiting distribution: 
            \begin{equation*}
                \sqrt{n}(\hat{\boldsymbol{\theta}} - \boldsymbol{\theta}_0) \to \mathcal{MN}(\textbf{0}, \boldsymbol{\Sigma}_{\theta}),
            \end{equation*}
        where $\boldsymbol{\Sigma}_{\theta} = \mathbb{E} \left( \frac{\partial \psi_i(\boldsymbol{\theta})}{ \partial \boldsymbol{\theta}} \biggr\rvert_{\boldsymbol{\theta}_0} \right)^{-1} 
                    \mathbb{E} [ \psi_i(\boldsymbol{\theta}_0)\psi_i(\boldsymbol{\theta}_0)^{\top} ]
                    \left\{ \mathbb{E} \left( \frac{\partial \psi_i(\boldsymbol{\theta})}{ \partial \boldsymbol{\theta}} \biggr\rvert_{\boldsymbol{\theta}_0} \right)^{-1} \right\}^{\top}$. \\                
        In order to define the estimating equations in terms of $\vect(\textbf{B})$, we need to introduce an equivalent representation of the model (\ref{equ:regmod2}), which is:
            \begin{equation*}
                \ObsRes = (\IdenMat_{T \times T} \otimes \ObsCov^{\top})\vect(\textbf{B}) + \boldsymbol{\epsilon}_i.
            \end{equation*}
        Based on this representation we can define the following equations, which together make up $\psi_i(\boldsymbol{\theta})$:
            \begin{equation}\label{equ:psi1}
                \psi_{1,i}(\boldsymbol{\theta}) = (\IdenMat_{T \times T} \otimes \ObsCov^{\top})\vect(\textbf{B}) + \frac{Z_i}{\tau(\ObsCov^{\top}\LogCoe)}(\textbf{y}_i - (\IdenMat_{T \times T} \otimes \ObsCov^{\top})\vect(\textbf{B})) - \boldsymbol{\mu}_{y}
            \end{equation}    
            \begin{equation}\label{equ:psi2}
                %\psi_{2,i}(\boldsymbol{\theta}) = Z_i \ObsCov \otimes (\textbf{y}_i - (\IdenMat_{T \times T} \otimes \ObsCov^{\top})\vect(\textbf{B}))
                \psi_{2,i}(\boldsymbol{\theta}) = (\ObsRes - (\IdenMat_{T \times T} \otimes \ObsCov^{\top})\vect(\textbf{B})) \otimes Z_i \ObsCov
            \end{equation}    
            \begin{equation}\label{equ:psi3}
                \psi_{3,i}(\boldsymbol{\theta}) = \frac{Z_i - \tau(\ObsCov^{\top} \LogCoe)}{\tau(\ObsCov^{\top} \LogCoe)(1-\tau(\ObsCov^{\top} \LogCoe))} \frac{\partial \tau(\ObsCov^{\top} \LogCoe)}{ \partial \LogCoe}.
            \end{equation}
        Note: (\ref{equ:psi1}), (\ref{equ:psi2}) and (\ref{equ:psi3}) all are vector valued functions and the dimensions are $T \times 1$, $Tp \times 1$, and $p \times 1$, respectively.
        %Let $G_n(\boldsymbol{\theta}) = n^{-1}\SumOverN \psi_i(\boldsymbol{\theta})$, where $\psi_i(\boldsymbol{\theta})$ is a vector valued equation containing equations (\ref{equ:psi1}) to (\ref{equ:psi3}), and the estimator $\hat{\boldsymbol{\theta}} = (\hat{\boldsymbol{\mu}}_y, \vect(\hat{\textbf{B}}), \hat{\LogCoe})^{\top}$ satisfies $G_n(\boldsymbol{\theta}) = \textbf{0}$.        
        Calculate the Jacobian of $\psi_i(\boldsymbol{\theta})$ and evaluate it at the true value $\boldsymbol{\theta}_0$.
        The first block of rows of the Jacobian is
            \begin{equation*}
                \frac{\partial \psi_{1,i}(\boldsymbol{\theta})}{ \partial \boldsymbol{\theta}^{\top} } \biggr\rvert_{\boldsymbol{\theta}_0}  
                = \left[ - \IdenMat_{T \times T}, \frac{\tau(\ObsCov^{\top} \LogCoe_0)- Z_i}{\tau(\ObsCov^{\top} \LogCoe_0 )} (\IdenMat_{T \times T} \otimes \ObsCov^{\top}), -Z_i e^{-\ObsCov^{\top} \LogCoe_0} \boldsymbol{\epsilon}_i \ObsCov^{\top} \right].
            \end{equation*}        
        The second block of rows is
            \begin{equation*}
                \frac{\partial \psi_{2,i}(\boldsymbol{\theta})}{ \partial \boldsymbol{\theta}^{\top} } \biggr\rvert_{\boldsymbol{\theta}_0} 
                = \left[ \textbf{0}_{Tp \times T},  Z_i \IdenMat_{T \times T} \otimes \textbf{x}_i\textbf{x}_i^{\top} , \textbf{0}_{p \times p} \right].
            \end{equation*}
        The last block of rows is
            \begin{equation*}
                \frac{\partial \psi_{3,i}(\boldsymbol{\theta})}{ \partial \boldsymbol{\theta}^{\top} } \biggr\rvert_{\boldsymbol{\theta}_0} 
                = \left[ \textbf{0}_{p \times T}, \textbf{0}_{p \times Tp}, \tau(\ObsCov^{\top} \LogCoe_0)(1-\tau(\ObsCov^{\top} \LogCoe_0)) \ObsCov\ObsCov^{\top}  \right].
            \end{equation*}
        
        If both working models are correctly specified, the expected value of the Jacobian is 
            \begin{equation*}
                \mathbb{E} \left( \frac{\partial \psi_i(\boldsymbol{\theta})}{ \partial \boldsymbol{\theta}} \biggr\rvert_{\boldsymbol{\theta}_0} \right) 
                    = \begin{bmatrix}
                    -\IdenMat_{T \times T} & \textbf{0}_{T \times Tp} & \textbf{0}_{T \times p}\\ 
                    \textbf{0}_{Tp \times T} & \IdenMat_{T \times T} \otimes \boldsymbol{\Pi}_{p \times p} & \textbf{0}_{p \times p}\\ 
                    \textbf{0}_{p \times T} & \textbf{0}_{p \times Tp} & \boldsymbol{\Phi}_{p \times p}
                    \end{bmatrix},
            \end{equation*}
        where $\boldsymbol{\Pi} = \mathbb{E}\left[ Z_i\textbf{x}_i\textbf{x}_i^{\top} \right]$ and $\boldsymbol{\Phi} = \mathbb{E}\left[\tau(\ObsCov^{\top} \LogCoe_0)(1-\tau(\ObsCov^{\top} \LogCoe_0)) \ObsCov\ObsCov^{\top}\right]$.
        To see why this is the case, note that if the missingness model \eqref{equ:ps} is correctly specified, the second block in the first row of the Jacobian simplifies as follows: 
            % \begin{equation*}
            %     \mathbb{E}\left(\frac{\tau(\ObsCov^{\top} \LogCoe_0)- Z_i}{\tau(\ObsCov^{\top} \LogCoe_0)} (\IdenMat_{T \times T} \otimes \ObsCov^{\top})\right)
            %     =\mathbb{E}_{\textbf{x}} \left( \mathbb{E} \left( \frac{\tau(\ObsCov^{\top} \LogCoe_0)- Z_i}{\tau(\ObsCov^{\top} \LogCoe_0)} (\IdenMat_{T \times T} \otimes \ObsCov^{\top}) \biggr\rvert \ObsCov \right) \right)
            %     = \textbf{0}.
            % \end{equation*}
            \begin{align*}
                & \mathbb{E}\left(\frac{\tau(\ObsCov^{\top} \LogCoe_0)- Z_i}{\tau(\ObsCov^{\top} \LogCoe_0)} (\IdenMat_{T \times T} \otimes \ObsCov^{\top})\right) \\
                &= \mathbb{E}_{\textbf{x}} \left( \mathbb{E} \left( \frac{\tau(\ObsCov^{\top} \LogCoe_0)- Z_i}{\tau(\ObsCov^{\top} \LogCoe_0)} (\IdenMat_{T \times T} \otimes \ObsCov^{\top}) \biggr\rvert \ObsCov \right) \right) = 0
            \end{align*}
        If the outcome model \eqref{equ:regmod2} is correctly specified, the third block in the first row of the Jacobian simplifies as follows:
            \begin{equation*}
                \mathbb{E} \left( -Z_i e^{-\ObsCov^{\top} \LogCoe_0} \boldsymbol{\epsilon}_i \ObsCov^{\top} \right) = \mathbb{E}\boldsymbol{\epsilon}_i \mathbb{E}\left( -Z_i e^{-\ObsCov^{\top} \LogCoe_0} \ObsCov^{\top} \right) = \textbf{0}. \\
            \end{equation*}    
        Next, calculate $\mathbb{E} [ \psi_i(\boldsymbol{\theta}_0)\psi_i(\boldsymbol{\theta}_0)^{\top} ] $.
        Still assuming both working models are correctly specified, we have: 
        %Therefore, we only need to specify $\mathbb{E} [ \psi_{j,i}(\boldsymbol{\theta}_0) \psi_{k,i}(\boldsymbol{\theta}_0)^{\top} ]$, where $j = 1,2,3$.
            \begin{align*}
                &\mathbb{E} [ \psi_{1,i}(\boldsymbol{\theta}_0) \psi_{1,i}(\boldsymbol{\theta}_0)^{\top} ] \\ 
                % &= \mathbb{E} \left[ \left( (\IdenMat_{T \times T} \otimes \ObsCov^{\top})\vect(\textbf{B}_0) - \boldsymbol{\mu}_y^0 + \frac{Z_i\epsilon_i}{\tau(\ObsCov^{\top}\LogCoe_0)} \right) 
                % \left( (\IdenMat_{T \times T} \otimes \ObsCov^{\top})\vect(\textbf{B}_0) - \boldsymbol{\mu}_y^0 + \frac{Z_i\epsilon_i}{\tau(\ObsCov^{\top}\LogCoe_0)} \right)^{\top} \right]  \\
                &= \Cov\left[ (\IdenMat_{T \times T} \otimes \ObsCov^{\top})\vect(\textbf{B}_0) \right] + \mathbb{E} \left[ \frac{Z_i}{\tau^2(\ObsCov^{\top}\LogCoe_0)} \boldsymbol{\epsilon}_i\boldsymbol{\epsilon}_i^{\top} \right] \\
                &= \textbf{B}^{\top}_0 \boldsymbol{\Sigma}_{x} \textbf{B}_0 + \mathbb{E} [\tau^{-1}(\ObsCov^{\top}\LogCoe_0)] \boldsymbol{\Sigma}_{\epsilon},
                %&= \mathbb{E}\left[ (\textbf{B}^{\top}_0\ObsCov - \mathbb{E}(\textbf{B}^{\top}_0\ObsCov | \ObsCov ))(\textbf{B}^{\top}_0\ObsCov - \mathbb{E}(\textbf{B}^{\top}_0\ObsCov | \ObsCov ))^{\top} \right]  
                %+ \mathbb{E} \left[ \frac{Z_i}{\tau^2(\ObsCov^{\top}\LogCoe_0)} \boldsymbol{\epsilon}_i\boldsymbol{\epsilon}_i^{\top} \right] \\
                %&= \mathbb{V}(\textbf{B}_0^{\top}\ObsCov) + \mathbb{E} [\tau^{-1}(\ObsCov^{\top}\LogCoe_0)] \boldsymbol{\Sigma}_{\epsilon} \\
                %&= \textbf{B}^{\top}_0 \boldsymbol{\Sigma}_{x} \textbf{B}_0 + \mathbb{E} [\tau^{-1}(\ObsCov^{\top}\LogCoe_0)] \boldsymbol{\Sigma}_{\epsilon}
            \end{align*}
        where 
            \begin{align*}
                \mathbb{E}\left[ \frac{1}{\tau(\ObsCov^{\top}\LogCoe_0)} \right] = \frac{1}{\pi}\left( 1 - \frac{\mathbb{V}(\tau(\ObsCov^{\top}\LogCoe_0))}{ [\mathbb{E}\tau(\ObsCov^{\top}\LogCoe_0)]^2 } \right) = \frac{1}{\pi} \left(1 - \frac{\mathbb{V}(\tau(\ObsCov^{\top}\LogCoe_0))}{\pi^2} \right),
            \end{align*}
            with $\pi=E(Z_i)$,
            and
            \begin{align*}
                \mathbb{V}(\tau(\ObsCov^{\top}\LogCoe_0)) &= \mathbb{E}_{\textbf{x}}[ \tau(\ObsCov^{\top}\LogCoe_0)^2 ] - \pi^2.
            \end{align*}
        Next,        
            \begin{align*}
                &
                \mathbb{E} [ \psi_{1,i}(\boldsymbol{\theta}_0) \psi_{2,i}(\boldsymbol{\theta}_0)^{\top} ] \\
                &= \mathbb{E} \left[ \left( (\IdenMat_{T \times T} \otimes \ObsCov^{\top})\vect(\textbf{B}_0) - \boldsymbol{\mu}_y^0 + \frac{Z_i}{\tau(\ObsCov^{\top}\LogCoe_0)}\boldsymbol{\epsilon}_i \right) \left( \boldsymbol{\epsilon}_i \otimes Z_i \ObsCov \right)^{\top} \right] \\
                &= \mathbb{E} \left[ \frac{Z_i}{\tau(\ObsCov^{\top}\LogCoe_0)} \boldsymbol{\epsilon}_i  \left( \boldsymbol{\epsilon}_i \otimes Z_i \ObsCov \right)^{\top} \right] 
                = \mathbb{E} \left[ \frac{Z_i}{\tau(\ObsCov^{\top}\LogCoe_0)}  ( \boldsymbol{\epsilon}_i \boldsymbol{\epsilon}_i^{\top} \otimes \ObsCov^{\top} ) \right] \\
                &= \boldsymbol{\Sigma}_{\epsilon} \otimes\boldsymbol{\mu}_x^{\top}
            \end{align*}
        and 
            \begin{align*}
                &\mathbb{E} [ \psi_{1,i}(\boldsymbol{\theta}_0) \psi_{3,i}(\boldsymbol{\theta}_0)^{\top} ] \\
                &= \mathbb{E} \left[ \left( (\IdenMat_{T \times T} \otimes \ObsCov^{\top})\vect(\textbf{B}_0) - \boldsymbol{\mu}_y^0 + \frac{Z_i}{\tau(\ObsCov^{\top}\LogCoe)}\boldsymbol{\epsilon}_i \right) (Z_i - \tau(\ObsCov^{\top}\LogCoe)) \ObsCov^{\top} \right]
                = \textbf{0}. \\
            \end{align*}
        Next,
            \begin{align*}
                &
                \mathbb{E} [ \psi_{2,i}(\boldsymbol{\theta}_0) \psi_{2,i}(\boldsymbol{\theta}_0)^{\top} ] \\
                &= \mathbb{E} \left[  (\boldsymbol{\epsilon}_i\otimes Z_i\ObsCov) (\boldsymbol{\epsilon}_i\otimes Z_i\ObsCov)^{\top} \right]
                = \mathbb{E}  \boldsymbol{\epsilon}_i\boldsymbol{\epsilon}_i^{\top} \otimes \left[ Z_i\ObsCov\ObsCov^{\top}  \right] \\
                &=\mathbb{E}\left[\boldsymbol{\epsilon}_i\boldsymbol{\epsilon}_i^{\top} \right] \otimes \mathbb{E} \left[Z_i \ObsCov\ObsCov^{\top}\right] 
                =   \boldsymbol{\Sigma}_{\epsilon} \otimes \boldsymbol{\Pi} 
            \end{align*}
        and 
            \begin{align*}
                \mathbb{E} [ \psi_{2,i}(\boldsymbol{\theta}_0) \psi_{3,i}(\boldsymbol{\theta}_0)^{\top} ] 
                =  \mathbb{E} \left[  ( \boldsymbol{\epsilon}_i \otimes Z_i\ObsCov ) (Z_i - \tau(\ObsCov^{\top} \LogCoe))\ObsCov^{\top} \right] = \textbf{0}.
            \end{align*}
        %Now, we are ready to find the asymptotic variance of the doubly robust estimator $\hat{\boldsymbol{\mu}}_y$.
        Since $\hat{\LogCoe}$ is a maximum likelihood estimator, the blocks in $\mathbb{E} [ \psi_{3,i}(\boldsymbol{\theta}_0)\psi_{3,i}(\boldsymbol{\theta}_0)^{\top} ] $ are identical to the expected value of the corresponding blocks in the Jacobian.
        According to the M-estimation approach, we have $\sqrt{n}( \hat{\boldsymbol{\theta}} -\boldsymbol{\theta} ) \overset{d}{\rightarrow} \mathcal{N}(\textbf{0}, \boldsymbol{\Sigma_\theta})$,
        where, by the sandwich estimator formula, the asymptotic variance of $\hat{\boldsymbol{\theta}}$ is
            \begin{align*}
                &  \mathbb{E} \left( \frac{\partial \psi_i(\boldsymbol{\theta})}{ \partial \boldsymbol{\theta}} \biggr\rvert_{\boldsymbol{\theta}_0} \right)^{-1} 
                    \mathbb{E} [ \psi_i(\boldsymbol{\theta}_0)\psi_i(\boldsymbol{\theta}_0)^{\top} ]
                    \left\{ \mathbb{E} \left( \frac{\partial \psi_i(\boldsymbol{\theta})}{ \partial \boldsymbol{\theta}} \biggr\rvert_{\boldsymbol{\theta}_0} \right)^{-1} \right\}^{\top} \\
                % &=  \begin{bmatrix}
                %         -\IdenMat & \textbf{0} & \textbf{0} \\ 
                %         \textbf{0}  & \IdenMat \otimes \boldsymbol{\Pi}^{-1} & \textbf{0}\\ 
                %         \textbf{0}  & \textbf{0} & \boldsymbol{\Phi}^{-1}
                %     \end{bmatrix}
                %     \begin{bmatrix}
                %         \textbf{B}^{\top}_0 \boldsymbol{\Sigma}_{x} \textbf{B}_0 + \mathbb{E} [\tau^{-1}(\ObsCov^{\top}\LogCoe_0)] \boldsymbol{\Sigma}_{\epsilon} &  \boldsymbol{\Sigma}_{\epsilon} \otimes \boldsymbol{\mu}_{\textbf{x}}^{\top}  & \textbf{0}\\ 
                %         \boldsymbol{\Sigma}_{\epsilon} \otimes \boldsymbol{\mu}_{\textbf{x}} &  \boldsymbol{\Sigma}_{\epsilon} \otimes \boldsymbol{\Pi}  & \textbf{0}\\ 
                %         \textbf{0} & \textbf{0} & \boldsymbol{\Phi}                                      
                %     \end{bmatrix}
                %     \begin{bmatrix}
                %         -\IdenMat & \textbf{0} & \textbf{0} \\ 
                %         \textbf{0}  & \IdenMat \otimes \boldsymbol{\Pi}^{-1} & \textbf{0}\\ 
                %         \textbf{0}  & \textbf{0} & \boldsymbol{\Phi}^{-1}
                %     \end{bmatrix} \\
                &=  \begin{bmatrix}
                        \textbf{B}^{\top}_0 \boldsymbol{\Sigma}_{x} \textbf{B}_0 + \mathbb{E} [\tau^{-1}(\ObsCov^{\top}\LogCoe_0)] \boldsymbol{\Sigma}_{\epsilon} &  -\boldsymbol{\Sigma}_{\epsilon} \otimes \boldsymbol{\mu}_x^{\top} \boldsymbol{\Pi}^{-1}  & \textbf{0} \\ 
                        \boldsymbol{\Sigma}_{\epsilon} \otimes \boldsymbol{\Pi}^{-1}\boldsymbol{\mu}_x & \boldsymbol{\Sigma}_{\epsilon} \otimes \boldsymbol{\Pi}^{-1} & \textbf{0} \\ 
                        \textbf{0} & \textbf{0} & \boldsymbol{\Phi}^{-1} 
                    \end{bmatrix}.
            \end{align*}
        Then this lemma is complete by the delta method.
        % At last, by the delta method, we have 
        %     $$
        %         \sqrt{n}(\hat{\boldsymbol{\mu}}_{DR} - \boldsymbol{\mu}_y^0) \overset{d}{\rightarrow} \mathcal{N}\left(\textbf{0}, \textbf{B}^{\top}_0 \boldsymbol{\Sigma}_{x} \textbf{B}_0 + \mathbb{E} [\tau(\ObsCov^{\top}\LogCoe_0)^{-1}] \boldsymbol{\Sigma}_{\epsilon}\right).
        %     $$
    \end{proof}

\end{document}
